\newcommand\version{v1}
\problemname{Kakor}
Ann Britt-Caroline har $N$ olika sorters kakor. Av kaksort $i$ har hon $A_i$ olika kakor. Nu undrar Ann hur många olika kakor hon har totalt sett.

\section*{Exempel}
I detta exempel har vi $N = 3$ olika kaksorter. Antalet kakor av varje sort är $3, 1, 5$.

\begin{figure}[h!]
  \centering
  \includegraphics[width=0.8\textwidth]{sample.png}
  \caption{Vi har tre kakor av första sorten, en av andra sorten, och fem av den tredje}
\end{figure}

Som ses på bilden ger detta totalt $9$ olika kakor.

\section*{Uppgift}
Du får givet $N$ och $A$. Skriv ett program som hjälper Ann att beräkna det totala antalet kakor. Du ska implementera funktionen \texttt{cookies(N, A)}:
\begin{itemize}
  \item \texttt{cookies(N, A)} - denna funktion kommer anropas exakt en gång av domaren.
  \begin{itemize}
    \item \texttt{N}: antalet olika kakor.
    \item \texttt{A}: en vektor av längd $N$. $A[i] (0 \le i < N) $ innehåller antalet kakor av sort $i$. $A[i]$ är alltid positiv.
    \item Funktionen ska returnera det totala antalet kakor som Anna har.
  \end{itemize}
\end{itemize}

\section*{Delpoäng}
Uppgiften består av ett antal grupper. Varje grupp ger ett visst antal poäng, och för att klara
gruppen måste du klara samtliga testfall i gruppen.

\begin{tabular}{|l|l|l|p{5cm}|}
  \hline
  \textbf{Grupp} & \textbf{Poäng} & \textbf{Gränser} \\ \hline
  1 & 27 & $1 \le N \le 1\,000$, $A[i] = 1$. \\ \hline
  2 & 34 & $1 \le N \le 1\,000$, $A[i] \le 100\,000$. \\ \hline
  3 & 39 & $1 \le N \le 100\,000$, $A[i] \le 100\,000$. \\ \hline
\end{tabular}

\section*{Indataformat}
Exempeldomaren läser indata i följande format:

\begin{itemize}
  \item rad 1: \texttt{N}
  \item rad 2: \texttt{A[0] A[1] ... A[N - 1]}
\end{itemize}

\section*{Utdataformat}
Exempeldomaren skriver ut returvärdet av \texttt{cookies(A, N)}.
